\begin{frame}{A hop is derived from the spaces, and a \texttt{Ref} names four residencies}
\centering
\begin{tikzpicture}[node distance=9mm,every node/.style={font=\scriptsize}]
  \node[box] (g) {\texttt{global}};
  \node[box,right=of g] (s) {\texttt{shared}};
  \node[box,right=of s] (r) {\texttt{register}};
  \draw[ar] (g) -- node[above,font=\tiny]{copy hop} node[below,font=\tiny,text=gray]{shared in \texttt{dsts}} (s);
  \draw[ar] (s) -- node[above,font=\tiny]{read hop} node[below,font=\tiny,text=gray]{shared in \texttt{srcs}} (r);
  \draw[ar,dashed] (g.north) to[out=40,in=140] node[above,font=\tiny]{one-hop} (r.north);
\end{tikzpicture}

\vspace{2pt}
{\scriptsize One \texttt{Move} class, no \texttt{Copy}/\texttt{Read} subclass --- so the analyses can walk
\emph{every} shared touch without a type switch, and "which movement is this?" is a
\bad{derived} property every new consumer must be told to get from \texttt{refs.py}.}

\vspace{5pt}
\begin{columns}[T]
\begin{column}{0.52\textwidth}
\scriptsize
\begin{tabular}{@{}ll@{}}
\toprule
\texttt{Ref} field group & residence \\
\midrule
\texttt{gen}, \texttt{gdelta}         & loop-carried shared \\
\texttt{abs\_gen}                     & pinned shared (peel, short arm) \\
\texttt{group}, \texttt{slot}, \texttt{reg\_ring} & register \\
\texttt{covers}                       & movement-quantum footprint \\
\bottomrule
\end{tabular}
\end{column}
\begin{column}{0.46\textwidth}
\scriptsize
\key{\texttt{gen} vs \texttt{abs\_gen}} is the peel/steady distinction made structural: a
prologue fill knows its buffer at compile time, a steady access does not.

\vspace{3pt}
\bad{"region" means two unrelated things} --- a \emph{storage} region (one piece of a split
LDS tile) and a \emph{placement} region (the legal window for a \texttt{Mark}).
They never share an expression, but they do share a file list.
\end{column}
\end{columns}
\end{frame}
